% Generated by scripts/assemble.py; corrections belong in source/corrections.json.
\examhead{108}{合肥工业大学 2025 年数学分析试卷}{106}
% p107-b17; source PDF page 107
\blockanchor{p107-b17}
\question{1}(10 分) 求极限  $\lim_{n\to\infty}\left(\frac{1}{n}-\frac{2}{n}+\frac{3}{n}+\cdots+\frac{(-1)^{n-1}n}{n}\right)$ .\par
% p107-b18; source PDF page 107
\blockanchor{p107-b18}
\question{2}(10 分) 求极限 $\lim_{x \to +\infty} e^{-x} \left(1 + \frac{1}{x}\right)^{x^2}$.\par
% p107-b19; source PDF page 107
\blockanchor{p107-b19}
\question{3}(15 分) 设 $f(x)$ 在 $[2, +\infty)$ 上连续, 且 $\lim_{n \to \infty} f(\sqrt{n}) = 0$. 证明: $\lim_{x \to +\infty} f(x)$ 存在当且仅当 $f(x)$ 在 $[2, +\infty)$ 上一致连续.\par
% p107-b20; source PDF page 107
\par\Needspace{6\baselineskip}
\blockanchor{p107-b20}
\question{4}(15 分) 设 $f(x)$ 在 $[0,1]$ 上可导, 且 $f(0) = 0, f(1) = 1$. 证明: 在 $(0,1)$ 内存在两个不同的点 $x_1, x_2$, 使得\par
\nopagebreak[4]
% p107-b21; source PDF page 107
\blockanchor{p107-b21}
\displayformula{\frac {1}{f ^ {\prime} (x _ {1})} + \frac {2024}{f ^ {\prime} (x _ {2})} = 2025.}
% p107-b22; source PDF page 107
\blockanchor{p107-b22}
\question{5}(15 分) 设函数 $f(x)$ 在 $[0,2]$ 上连续可导, 且 $f(0) = f(2) = 1$, $|f'(x)| \leq 1$, 证明: $\int_0^2 f(x) \, \mathrm{d}x > 1$.\par
% p107-b23; source PDF page 107
\blockanchor{p107-b23}
\question{6}(10 分) 讨论级数 $\sum_{n=1}^{\infty} \frac{(-1)^{n-1}}{n} \frac{a}{1 + a^n} (a > 0)$ 的绝对收敛和条件收敛性.\par
% p107-b24; source PDF page 107
\blockanchor{p107-b24}
\question{7}(15 分) 设 $\{f_n(x)\}$ 为 $[a, b]$ 上的 Riemann 可积函数列, 且函数列 $\{f_n(x)\}$ 在 $[a, b]$ 上一致收敛于 $f(x)$. 证明: $f(x)$ 在 $[a, b]$ 上 Riemann 可积, 且 $\lim_{n \to \infty} \int_{a}^{b} f_n(x) \mathrm{d}x = \int_{a}^{b} f(x) \mathrm{d}x$.\par
% p108-b00; source PDF page 108
\par\Needspace{6\baselineskip}
\blockanchor{p108-b00}
\question{8}(15 分) 抛物面 $x^{2} + y^{2} = z$ 被平面 $x + y + z = 1$ 截成一个椭圆, 求这个椭圆到原点的最长与最短距离.
\question{9}(15 分) 计算曲面积分\par
\nopagebreak[4]
% p108-b01; source PDF page 108
\blockanchor{p108-b01}
\displayformula{\iint_ {S} y z \mathrm{d} y \mathrm{d} z + (x ^ {2} + z ^ {2}) y \mathrm{d} z \mathrm{d} x + x y \mathrm{d} x \mathrm{d} y.}
% p108-b02; source PDF page 108
\blockanchor{p108-b02}
其中 $S$ 为曲面 $y = 4 - (x^{2} + y^{2})(y > 0)$, 在 $xoz$ 面右侧部分外侧.\par
% p108-b03; source PDF page 108
\par\Needspace{6\baselineskip}
\blockanchor{p108-b03}
\question{10}(15 分) 设 $z = f(x - y, x + y) + g(x + ky)$, $f$ 具有二阶连续偏导数, $g$ 具有二阶连续导数且 $g''(x)$ 不恒为 0. 如果\par
\nopagebreak[4]
% p108-b04; source PDF page 108
\blockanchor{p108-b04}
\displayformula{\frac {\partial^ {2} z}{\partial x ^ {2}} + 2 \frac {\partial^ {2} z}{\partial x \partial y} + \frac {\partial^ {2} z}{\partial y ^ {2}} = 4 f _ {22} ^ {\prime \prime}.}
% p108-b05; source PDF page 108
\blockanchor{p108-b05}
求常数 $k$ 的值.\par
% p108-b06; source PDF page 108
\blockanchor{p108-b06}
\question{11}(15 分) 求函数 $I(\alpha) = \int_{1}^{+\infty} \frac{\sin x \cos x}{x^{\alpha}} \, \mathrm{d}x$ 的定义域, 并研究其在定义域上的连续性.\par
